% !PRINTED 158
\weberpage{158}

Die Berechnung der Resultanten ist meist sehr weitläufig. Wir wollen nur das einfachste Beispiel der Resultante zweier quadratischer Functionen anführen:
\[
f(x) = a_0 x^2 + a_1 x + a_2
\]
\[
\varphi(x) = b_0 x^2 + b_1 x + b_2.
\]
Man kann hier direct das Product
\[
R = a_0^2 (b_0 \alpha_1^2 + b_1 \alpha_1 + b_2) (b_0 \alpha_2^2 + b_1 \alpha_2 + b_2)
\]
berechnen, und erhält, wenn man
\[
- a_0 (\alpha_1 + \alpha_2) = a_1, \qquad a_0 \alpha_1 \alpha_2 = a_2
\]
setzt,
\weq{11}{
R = a_0^2 b_2^2 + a_2^2 b_0^2 - b_0 b_1 a_1 a_2 - a_0 a_1 b_1 b_2
+ b_0 b_2 a_1^2 + a_0 a_2 b_1^2 - 2 a_0 b_0 a_2 b_2,
}
oder was damit identisch ist
\weq{12}{
R = (a_0 b_2 - a_2 b_0)^2 - (a_0 b_1 - a_1 b_0) (a_1 b_2 - a_2 b_1).
}
Die einzelnen Glieder der Resultante der zwei Functionen $f(x)$ und $\varphi(x)$ vom $n^{\mathrm{ten}}$ und $m^{\mathrm{ten}}$ Grade haben, von einem numerischen (ganzzahligen) Factor abgesehen, die Gestalt
\weq{13}{
a_0^{\nu_0} a_1^{\nu_1} \ldots a_n^{\nu_n} b_0^{\mu_0} b_1^{\mu_1} \ldots b_m^{\mu_m},
}
worin, wie aus den oben bestimmten Graden hervorgeht,
% [source-original mathematical slip, confirmed 2026-08-25 against scan p. 183 / printed p. 158]
% The printed right-hand sides are reversed: since R is homogeneous of degree m in the
% a-coefficients and degree n in the b-coefficients (as stated on printed p. 156), the
% mathematically consistent relations would be sum(nu_i)=m and sum(mu_i)=n. Preserved verbatim.
\weq{14}{
\begin{aligned}
\nu_0 + \nu_1 + \cdots + \nu_n &= n \\
\mu_0 + \mu_1 + \cdots + \mu_m &= m.
\end{aligned}
}
Wir können aber noch eine andere Relation zwischen diesen Exponenten angeben.

Da nämlich [nach (6)] $R$ eine homogene Function $n\,m^{\mathrm{ten}}$ Grades von den $n+m$ Variablen $\alpha, \beta$ ist, da ferner $a_0$ vom nullten, $a_1$ vom ersten, $a_2$ vom zweiten etc. Grade in den $\alpha$, allgemein $a_k$ und $b_k$ vom $k^{\mathrm{ten}}$ Grade in diesen Variablen sind, so folgt
\weq{15}{
\begin{aligned}
\nu_1 + 2 \nu_2 + 3 \nu_3 + \cdots + n \nu_n \\
{}+ \mu_1 + 2 \mu_2 + 3 \mu_3 + \cdots + m \mu_m &= n m,
\end{aligned}
}
eine Relation, aus der sich ein wichtiger Schluss ziehen lässt, den wir im folgenden Paragraphen etwas ausführlicher besprechen wollen.
